%% kapa: exam=1 round=37 year=2026 period=2 form=A subject=appraisal_law subject_ko=감정평가관계법규 seq=21 number=21
%% source: Q-Net 감정평가사 기출문제 (cst00309) artlSeq=5258924 file=raw/1008/5258924/2026년 제37회 감정평가사 1차 2교시.pdf
%% conversion: kordoc 4.24.51 -> tools/gen1.py (auto); review=proofread-sonnet; flags=none
\begin{question}{21}
국유재산법령상 일반재산에 관한 설명으로 옳은 것은?
\begin{choices}
  \item 총괄청은 일반재산을 보존용재산으로 전환하여 관리할 수 없다.
  \item 일반재산은 ｢민법｣ 제245조에도 불구하고 시효취득의 대상이 되지 아니한다.
  \item 재판상 화해에 따라 일반재산에 사권을 설정할 수 있다.
  \item 일반재산은 조림 사업을 시행하기 위하여 그 사업의 완성을 조건으로 양여를 예약할 수 있고, 그 예약기간은 계약일부터 5년 이내로 정하여야 한다.
  \item 총괄청은 용도를 폐지함으로써 일반재산으로 된 국유재산에 대하여 필요하다고 인정하는 경우에는 그 처리방법을 지정해야 하며, 이를 직접 처리해서는 안 된다.
\end{choices}
\end{question}
\answer{3}
